% ch10_conclusion.tex — Dissertation Chapter 10: Conclusion
% Blair Dupre, PhD Dissertation, University of North Dakota

\chapter{Conclusion}
\label{ch:conclusion}

This dissertation has presented the first end-to-end NSD-ISS computational pipeline integrating classification, graph-informed imputation, benchmarked transition-timing prediction, conformalized uncertainty quantification, temporal deployment validation, and unified research-setting clinical decision support for Parkinson's disease. Across six papers, the GIMAN framework addresses a gap confirmed by systematic review~\cite{page2021}---zero computational tools targeting NSD-ISS staging~\cite{simuni2024} existed prior to this work---and provides a complete pipeline from raw clinical data to uncertainty-aware clinical decision support.

Paper~1 (Chapter~\ref{ch:paper1}) established that NSD-ISS biological stages can be predicted from multimodal clinical features with high discrimination. CatBoost achieved AUC 0.981 for binary NSD-positive detection and 0.944 for three-class staging across 46 features and 10 modalities, with conformal prediction sets~\cite{vovk2022} guaranteeing at least 90\% coverage of the true stage. External validation across four independent cohorts (PPMI, BioFIND, PDBP, HBS) revealed a critical domain shift: PPMI's inclusion of healthy controls confounds binary classification on PD-only external cohorts, while NSD-positive sub-staging generalizes robustly (AUC 0.900 with clinical features alone). This finding has direct implications for how future NSD-ISS prediction studies should define their training populations.

Paper~2 (Chapter~\ref{ch:paper2}) demonstrated that graph-informed imputation conditioned on biological stage produces superior downstream clinical utility despite a paradoxical relationship with aggregate imputation error. GIMIN achieved RMSE 107.7---22\% lower than MissForest~\cite{stekhoven2012}, the strongest of eight classical and deep learning baselines across a comprehensive 12-model benchmark. The imputation-utility paradox, wherein stage conditioning slightly worsens aggregate RMSE but improves downstream balanced accuracy by 2.9--3.1\% on clinically relevant targets, reveals that imputation should be optimized for the end task rather than for intermediate reconstruction fidelity. This insight applies broadly to any clinical prediction problem with class imbalance and missing data.

Paper~3 (Chapter~\ref{ch:paper3}) presented the first benchmarked comparison of multi-state Markov, Dynamic-DeepHit, and Graph-Informed Digital Twin models for NSD-ISS stage transition timing. Using 2,859 observed transitions from 1,900 longitudinally staged patients, Dynamic-DeepHit~\cite{lee2018} achieved C-td 0.926, correctly ranking 92.6\% of patient pairs by their transition timing. The Graph-Informed Digital Twin matched this discrimination (C-td 0.926) and showed numerically lower fold-to-fold standard deviation (0.013 versus 0.018), though this difference did not reach statistical significance (paired $t$-test $p{=}0.345$). The longitudinal analysis revealed that 39.1\% of observed transitions are backward, demonstrating that NSD-ISS stages are not strictly progressive and that medications can reverse biological stage assignments~\cite{espay2025,simuni2025}---a finding with important implications for clinical trial design and patient communication.

Paper~4 (Chapter~\ref{ch:paper4}) developed the first conformalized survival analysis framework for competing-risks cumulative incidence functions in the NSD-ISS setting. Custom IPCW-weighted conformal CIF bands achieved 91.1\% marginal coverage at 95\% confidence level, with band widths 2.6 times narrower than naive conformal approaches. Calibration analysis confirmed ECE below 0.005 at all clinical horizons, and bootstrap interaction tests with Benjamini-Hochberg FDR correction~\cite{benjamini1995} demonstrated equitable performance across sex, age, and genetic subgroups with no significant model-by-subgroup interactions. These results establish that survival predictions for NSD-ISS transitions can be made both accurate and honestly calibrated.

Paper~5 (Chapter~\ref{ch:paper5}) evaluated deployment readiness through expanding-window temporal validation, training on historically enrolled patients and testing on genuinely future ones. The models maintained strong discrimination (C-td above 0.85) through three realistic temporal windows, with temporal degradation of 9.3\% for DeepHit and 9.0\% for Graph-DT under the extreme stress-test scenario. Multivariate covariate shift detection identified which clinical features drift across enrollment cohorts, providing the infrastructure for proactive retraining triggers in a deployed system.

Paper~6 (Chapter~\ref{ch:paper6}) integrated all preceding components into a unified clinical decision support pipeline demonstrated through five patient case studies. For any patient, the system provides NSD-ISS stage classification with class probabilities, cause-specific cumulative incidence functions from both survival models, and conformal prediction bands with calibrated coverage---all within a latency of under two seconds. The pipeline demonstrates that the technical components developed across Papers~1--5 can be composed into a coherent system suitable for clinical workflow integration.

The clinical implications of this work extend across three stakeholder groups. For neurologists and movement disorder specialists, the framework provides the first personalized prognostic tool for NSD-ISS biological staging, replacing population-average statistics with individual-level predictions grounded in each patient's unique multimodal profile. The ``patients like you'' paradigm enabled by the graph-informed approach mirrors how clinicians naturally reason by analogy, and the conformal prediction bands communicate model uncertainty in a form suitable for shared decision-making with patients. For pharmaceutical and biotechnology companies developing disease-modifying therapies, the transition timing predictions offer a mechanism for clinical trial enrichment---identifying patients most likely to progress within a trial timeframe, potentially reducing required sample sizes by 30--50\%~\cite{temple2010} and saving tens of millions of dollars per trial~\cite{wouters2020}. For patients and caregivers, the quantitative prognostic information enables informed life planning, and the finding that 39.1\% of transitions are backward provides evidence-based hope that current treatments can meaningfully alter the biological trajectory of the disease.

This dissertation establishes the computational infrastructure for NSD-ISS-based precision neurology. The framework is designed to be general: the same pipeline architecture---graph-informed imputation, tree-based classification, deep survival analysis with graph fusion, conformal uncertainty quantification, temporal deployment validation, and unified clinical reporting---can be retargeted to any future biological staging system by changing the target labels. Post-dissertation work (Phase 2 of the mechanistic digital twin extension) builds on this foundation with a narrow, distinct contribution: the first per-patient practical identifiability analysis of the $\alpha$-synuclein oligomer $\to$ dopaminergic neuron loss coupling inferred from longitudinal serial DaT-SPECT imaging, principally reporting the toxicity flux composite $T_{\text{tox}} = \alpha_{\text{tox}} \cdot k_n \cdot M_{ss}^2 / (k_{\text{conv}} + k_{\text{clear},O})$ rather than individual rate constants that are practically non-identifiable under slow-fast timescale collapse~\cite{villaverde2016,raue2009,gutenkunst2007,transtrum2015}. This extension is explicitly distinct from, and complementary to, parallel quantitative systems pharmacology and digital twin efforts including the Critical Path for Parkinson's consortium framework~\cite{denaro2024} and the clinical-score latent state-space twin of Hemedan et al.~\cite{hemedan2026}, neither of which reports mechanistic compartmental biochemistry with practical identifiability diagnostics. As the NSD-ISS framework matures and as mechanistic understanding of alpha-synuclein pathology deepens, the correlational predictions of GIMAN will serve as the foundation upon which truly interventional mechanistic digital twins are built---models that can simulate not just what will happen, but what could happen under the treatments we choose to give. The ultimate vision is a clinical tool where a neurologist enters a patient's data and receives not only a prognosis but a personalized treatment recommendation, grounded in calibrated uncertainty, validated across diverse populations, and continuously updated as new data arrives. This dissertation takes the first comprehensive step toward that vision.
